% ========== 第七章：总结与展望 ==========
\chapter{总结与展望}

\section{工作总结}

本文围绕"基于扩散模型的多视角相机镜头控制图像生成"这一核心任务，提出了一套以 InstructPix2Pix 为底模、以 LoRA 为轻量微调手段的多模态条件扩散方案。主要工作可以概括为以下四个方面：

\textbf{第一，方法层面}，提出了结构图视角变换的新范式：将抽象的相机控制问题转化为扩散模型可理解的结构图变换问题，以深度图/轮廓图作为中间条件表示，复用 InstructPix2Pix 原生的 8 通道 UNet 实现条件 latent 注入，无需从零训练条件权重。设计了 RotationEncoder（sin/cos 编码 + MLP），将三维旋转向量转换为可与 CLIP 文本嵌入拼接的伪文本 token，使相机位姿信息进入 UNet cross-attention 路径。

\textbf{第二，数据层面}，搭建了基于 Blender 的全自动化多视角合成数据管线，覆盖场景搭建、分层网格相机轨迹采样、批量渲染、Depth Anything V2 深度图/Canny 轮廓图后处理、以及按场景拆分的 train/val 划分全流程。管线支持内置几何体降级与外部 3D 模型扩展，可通过命令行参数灵活调节场景数、视角数、分辨率、渲染引擎等。

\textbf{第三，训练层面}，设计了相邻视角配对策略，将训练任务从"任意角度旋转"收窄为"小角度相邻旋转变换"，降低了单步学习难度，提升了 loss 收敛稳定性。结合双向训练消除方向偏置，并通过多步迭代推理将大角度旋转分解为串联的小步长序列，使模型仅需小角度能力即可覆盖任意旋转幅度。

\textbf{第四，工程层面}，提供了完整的训练—验证—推理—部署链路：LoRA 训练脚本支持断点续训、显存优化、混合精度与梯度累积；自动化收敛验证脚本无需 GPU 即可检测 loss 下降趋势与稳定性；Gradio 推理界面支持双分支独立生成、多步迭代可视化、以及 LoRA 缺失时自动退化为底模原生图生图；ComfyUI 工作流提供可视化节点部署方案。

\section{未来工作展望}

基于本文的工作基础与第~\ref{sec:limitations} 节讨论的局限性，未来可在以下方向继续深入：

\begin{enumerate}[label=(\arabic*)]
    \item \textbf{显式 6DoF 相机位姿编码}：将当前仅包含旋转向量的条件扩展为完整的 6DoF 位姿表示（位置 + 旋转），例如采用 Plücker 射线或相机外参矩阵作为条件，使模型具备真正的三维相机控制能力。
    
    \item \textbf{端到端 RGB 图像生成}：将当前"结构图变换 + ControlNet RGB 生成"的两阶段管线整合为端到端训练，直接从参考 RGB 图与旋转向量生成目标视角的 RGB 图像。这需要在更大规模的数据集上训练，并可能需要引入感知损失（如 LPIPS）或对抗损失以提升生成质量。
    
    \item \textbf{多视角一致性约束}：引入跨视角 attention 机制或几何一致性损失（如在 3DGS/NeRF 隐式表示空间中施加光度一致性约束），使多个目标视角的生成结果之间保持结构、身份、材质与光照的一致性。
    
    \item \textbf{更大规模与更多样化的数据集}：使用更多样的 3D 模型资产（如 Objaverse、ShapeNet），添加更丰富的背景、材质与光照变化，缩小合成数据与真实图像的 domain gap。同时可考虑在真实多视角数据集（如 CO3D、MVImgNet）上进行微调。
    
    \item \textbf{roll 角显式训练}：在 Blender 渲染阶段显式添加相机 roll 角变化，或在推理时通过对生成结果做后处理旋转变换来获得 roll 角标签，使模型具备完整的 3 自由度旋转控制。
    
    \item \textbf{连续镜头轨迹生成}：利用多步迭代策略的自然扩展，实现平滑的连续镜头轨迹（camera path）生成，应用于动画分镜、虚拟场景漫游等场景。
    
    \item \textbf{真实场景微调与评估}：在真实拍摄的多视角图像上对模型进行微调，并以人工评估或用户调研的方式验证模型在实际应用中的镜头控制准确性与生成质量。
\end{enumerate}

综上所述，本文构建的原型系统验证了"通过中间结构条件让扩散模型服从镜头控制"这一核心思路的可行性，为后续的多视角一致生成、连续镜头生成与显式三维感知的扩散模型研究提供了可复现的基础与参照。
